%!TEX root = ./main.tex

\section[The conflict]{Why One Signal Cannot Be Great at Both}

% SECTION TIMING: 24:00--32:00 (3 frames). Question 1 of the roadmap frame.
% TEACHING NOTE: the two tables are one table in the source talk. They are split here
% because ten rows at a readable size do not fit on one frame. Read down a column, then
% across one row. Do not read all twenty cells aloud.

\begin{frame}{They want opposite things (1/2)}
  \vspace{0.2em}
  \scriptsize
  \setlength{\tabcolsep}{4pt}
  \renewcommand{\arraystretch}{1.35}
  \begin{tabularx}{\textwidth}{>{\bfseries\raggedright\arraybackslash}p{1.55cm}
      >{\raggedright\arraybackslash}X >{\raggedright\arraybackslash}X}
    \toprule
    & \textcolor{coreorange}{\textbf{Sensing wants}} & \textcolor{ranblue}{\textbf{Communication wants}} \\
    \midrule
    Spectrum   & Wide bandwidth, for resolution & Efficient bandwidth, for rate \\
    Signal     & Deterministic waveforms & Random, information-bearing symbols \\
    Power      & High transmit power on the target & Power efficiency per bit \\
    Time       & Continuous illumination & Low-latency bursts \\
    Space      & Wide or scanning beams & Narrow beams, one per user \\
    \bottomrule
  \end{tabularx}

  \vspace{0.6em}
  \takeaway{Randomness is what carries information --- and randomness ruins a clean echo.\\That single sentence is the whole conflict.}
\end{frame}

\begin{frame}{They want opposite things (2/2)}
  \vspace{0.2em}
  \scriptsize
  \setlength{\tabcolsep}{4pt}
  \renewcommand{\arraystretch}{1.35}
  \begin{tabularx}{\textwidth}{>{\bfseries\raggedright\arraybackslash}p{1.55cm}
      >{\raggedright\arraybackslash}X >{\raggedright\arraybackslash}X}
    \toprule
    & \textcolor{coreorange}{\textbf{Sensing wants}} & \textcolor{ranblue}{\textbf{Communication wants}} \\
    \midrule
    Interference & A clean echo, no clutter & A clean constellation, no radar leakage \\
    Domain     & Fine steps in distance and speed & Fine slots in time and frequency \\
    Hardware   & Precise timing, high-power amplifiers & Spectral efficiency, many users \\
    Latency    & Low-latency tracking feedback & Tolerates delay off the critical path \\
    Regulation & Frequency allocation rules & Emission and efficiency standards \\
    \bottomrule
  \end{tabularx}

  \vspace{0.6em}
  \qa{Is there a setting where both are optimal?}
     {No. There is a \textbf{Pareto frontier}: past a point, every improvement in
      one is paid for out of the other. Design work is choosing where to sit on
      it, not escaping it.}
\end{frame}

% ACCURACY NOTE: the horizontal axis is sensing ERROR, so the good corner of this plot
% is up and to the LEFT -- and the frontier is precisely what keeps that corner out of
% reach. The curve therefore RISES to the right: allowing a larger sensing error buys a
% higher rate. Drawing it falling would put the sensing-first point above AND left of the
% comms-first point, i.e. better on both axes, which is domination, not a tradeoff.
% Students read every curve as "up and to the right is better" and will get the axis
% backwards unless you say it. (The axis is the CRLB in the source paper; the acronym is
% kept out of the main deck for this audience.)
\begin{frame}{The frontier, drawn once}
  \vspace{0.3em}
  \centering
  \begin{tikzpicture}[scale=0.95,transform shape]
    \draw[-{Latex[length=2.2mm]},softgray] (0,0) -- (6.6,0)
      node[anchor=north east,font=\scriptsize] {sensing error $\rightarrow$ worse};
    \draw[-{Latex[length=2.2mm]},softgray] (0,0) -- (0,3.6)
      node[anchor=south,font=\scriptsize] {achievable rate};
    \draw[very thick,draw=ranblue] (0.7,0.9) .. controls (2.0,2.3) and (3.2,2.85) .. (5.6,3.15);
    \node[font=\scriptsize,text=ranblue,anchor=north west] at (3.15,2.45) {attainable frontier};
    \fill[softgreen] (1.68,1.80) circle (2.6pt);
    \node[font=\scriptsize,text=softgreen,anchor=north west] at (1.82,1.68) {sensing-first};
    \fill[coreorange] (4.60,3.00) circle (2.6pt);
    \node[font=\scriptsize,text=coreorange,anchor=west] at (4.75,3.30) {comms-first};
    \node[font=\scriptsize,text=softgray] at (1.15,2.95) {infeasible};
    \node[font=\scriptsize,text=softgray] at (4.45,1.10) {reachable, but wasteful};
  \end{tikzpicture}

  \vspace{0.6em}
  \glossline{Horizontal axis = sensing error: how far off the range and speed estimates
    are. \textbf{Lower is better}, so the good corner is up and to the left --- and the
    curve is exactly what keeps you out of it: walking left buys accuracy with rate.}

  \vspace{0.2em}
  \takeaway{Every result in the rest of this lecture\\is a claim about one point on a curve like this.}
\end{frame}
